\documentclass[11pt,a4paper,fontset=fandol]{ctexart}
\usepackage[a4paper,margin=1.78cm,headheight=15pt]{geometry}
\usepackage{amsmath,amssymb,mathtools,bm}
\usepackage{booktabs,longtable,tabularx,array}
\usepackage{enumitem}
\usepackage[most]{tcolorbox}
\usepackage{tikz}
\usetikzlibrary{arrows.meta,positioning,calc}
\usepackage{xcolor}
\usepackage{fancyhdr}
\usepackage{hyperref}
\usepackage{bookmark}
\usepackage{microtype}

\newcolumntype{Y}{>{\raggedright\arraybackslash}X}
\newcolumntype{L}[1]{>{\raggedright\arraybackslash}p{#1}}
\setlength{\parindent}{2em}
\setlength{\parskip}{0.34em}
\setlength{\emergencystretch}{3em}
\renewcommand{\arraystretch}{1.22}
\setlength{\tabcolsep}{3pt}
\setlist[itemize]{itemsep=0.18em,topsep=0.35em,leftmargin=2em}
\setlist[enumerate]{itemsep=0.18em,topsep=0.35em,leftmargin=2em}

\definecolor{deep}{RGB}{42,58,73}
\definecolor{soft}{RGB}{245,247,249}
\definecolor{accent}{RGB}{104,76,55}
\definecolor{greenish}{RGB}{57,92,76}
\definecolor{warn}{RGB}{136,68,63}
\definecolor{purple}{RGB}{87,72,116}
\hypersetup{colorlinks=true,linkcolor=deep,urlcolor=accent,pdftitle={定向积分与向量场}}
\pagestyle{fancy}
\fancyhf{}
\lhead{\small\color{deep}\textbf{数学一 · 高等数学 Topic 09}}
\rhead{\small\color{deep}曲线 · 曲面 · 定向 · 场}
\cfoot{\small\thepage}
\renewcommand{\headrulewidth}{0.4pt}
\newtcolorbox{corebox}[1]{breakable,colback=soft,colframe=deep,title=\textbf{#1},boxrule=0.8pt,arc=2mm}
\newtcolorbox{methodbox}[1]{breakable,colback=white,colframe=greenish,title=\textbf{#1},boxrule=0.7pt,arc=1.5mm}
\newtcolorbox{warnbox}[1]{breakable,colback=white,colframe=warn,title=\textbf{#1},boxrule=0.7pt,arc=1.5mm}
\newtcolorbox{examplebox}[1]{breakable,colback=white,colframe=purple,title=\textbf{#1},boxrule=0.7pt,arc=1.5mm}
\newcommand{\kw}[1]{\textbf{\color{deep}#1}}

\begin{document}
\begin{center}
{\LARGE\textbf{定向积分与向量场}}\\[0.4em]
{\large 从“微元分类”到“内部微分量与边界积分的转换”}
\end{center}
\vspace{0.5em}

\begin{quote}
本册只追问一个根本问题：\textbf{曲线或曲面上的局部贡献究竟由什么场、什么几何微元和什么方向共同决定；当边界与内部满足正则性时，怎样用 Green、Stokes、Gauss 在“复杂边界”与“内部微分量”之间换表示，从而降低总计算复杂度？} 这些定理并不保证降维：曲线常换成面积分、曲面常换成体积分；真正的收益是把难处理的方向、分片和参数化换成更简单的旋度、散度与内部区域。看见积分符号后，第一动作不是代参数，而是辨认场的类型、微元的载体和定向协议。
\end{quote}

\begin{center}\rule{0.5\linewidth}{0.5pt}\end{center}

\setcounter{section}{-1}
\section{本册定位}

本册是高等数学 Subject Atlas 下的 Topic09。Topic06 提供空间对象与参数/约束表示，Topic08 提供无向面积体积累积；本册把曲线、曲面、方向和向量场组织成定向积分与边界转换系统。

\subsection{Own}
\begin{itemize}
  \item 第一型曲线/曲面积分与长度/面积微元；
  \item 第二型曲线积分的切向位移与第二型曲面积分的法向面积；
  \item 参数化、投影、直接代入、对称性和物理量建模；
  \item 梯度场、保守场、路径无关和穿孔区域反例；
  \item Green、Stokes、Gauss 三大定理的对象、方向、光滑性与奇点资格；
  \item 补线、补面、挖洞和“换面不换边界”的统一决策接口。
\end{itemize}

\subsection{Uses / Stop Boundary}
\begin{itemize}
  \item Topic06 提供曲线/曲面对象、参数域与法向表示；
  \item Topic08 提供二重/三重积分、投影区域、坐标换元与无向微元；
  \item Topic07 提供梯度、方向导数和局部可微语言；
  \item B00 拥有内积/正交/投影的跨学科统一解释；
  \item 一般微分形式、流形、拓扑同调和偏微分方程理论不进入数学一主干；
  \item 单一定向积分/向量场问题族的触发信号、第一动作、检查点和退出条件进入同目录训练 Markdown；跨多个训练专题的稳定控制才进入《高等数学做题规则》。
\end{itemize}

\section{母模型：场 × 载体 × 定向}

\[
\boxed{\begin{gathered}
\text{Field Type / Geometric Carrier}
\to\text{Orientation / Parametrization}
\to\text{Microelement}\\
\to\text{Direct or Boundary Theorem}
\to\text{Domain / Singularity Audit}
\to\text{Invariant Check}
\end{gathered}}
\]

\begin{enumerate}
  \item \kw{Field Type / Carrier}：标量场配长度/面积，向量场配切向位移/法向面积；
  \item \kw{Orientation / Parametrization}：确认曲线方向、曲面法向、边界右手关系；
  \item \kw{Microelement}：分别写 $ds,dS,d\mathbf r,d\mathbf S$，不可混用；
  \item \kw{Direct or Boundary Theorem}：比较参数化、投影、保守场和三大定理；
  \item \kw{Domain / Singularity Audit}：检查闭合、单连通、光滑性和内部奇点；
  \item \kw{Invariant Check}：反向、换面、量纲、对称和简单场值必须自洽。
\end{enumerate}

\begin{corebox}{四类积分不是四套孤立公式}
\[
\boxed{\begin{array}{ccc}
\toprule
&\text{曲线}&\text{曲面}\\ \midrule
\text{标量场}&\displaystyle\int_L f\,ds&\displaystyle\iint_\Sigma f\,dS\\[5pt]
\text{向量场}&\displaystyle\int_L\mathbf F\cdot d\mathbf r&\displaystyle\iint_\Sigma\mathbf F\cdot d\mathbf S\\
\bottomrule
\end{array}}
\]
第一型积分只累积无向密度；第二型积分从向量场提取切向或法向分量，因此反向会变号。
\end{corebox}

\begin{center}
\begin{tikzpicture}[
  node distance=0.58cm and 0.58cm,
  box/.style={draw=deep,rounded corners=1.5mm,minimum width=3.1cm,minimum height=1.0cm,align=center,fill=soft,font=\small},
  arr/.style={-{Latex[length=2mm]},thick,draw=deep}
]
\node[box] (a) {场与对象\\标量/向量 · 线/面};
\node[box,right=of a] (b) {定向与参数\\切向/法向/边界};
\node[box,right=of b] (c) {微元\\$ds,dS,d\mathbf r,d\mathbf S$};
\node[box,below=of c] (d) {直接计算\\代入/投影/对称};
\node[box,left=of d] (e) {边界转换\\Green/Stokes/Gauss};
\node[box,left=of e] (f) {资格与复核\\奇点/拓扑/符号};
\draw[arr] (a)--(b); \draw[arr] (b)--(c); \draw[arr] (c)--(d);
\draw[arr] (d)--(e); \draw[arr] (e)--(f);
\end{tikzpicture}
\end{center}

\section{第一型：标量密度乘无向几何微元}

若曲线参数化为 $\mathbf r(t)$，则
\[
ds=\lVert\mathbf r'(t)\rVert\,dt,
\qquad
\int_L f\,ds=\int_\alpha^\beta f(\mathbf r(t))\lVert\mathbf r'(t)\rVert\,dt.
\]
弧长因子非负，所以反向参数化不改变第一型积分。

若曲面参数化为 $\mathbf r(u,v)$，则
\[
dS=\lVert\mathbf r_u\times\mathbf r_v\rVert\,du\,dv,
\qquad
\iint_\Sigma f\,dS
=\iint_D f(\mathbf r)\lVert\mathbf r_u\times\mathbf r_v\rVert\,du\,dv.
\]
显式曲面 $z=z(x,y)$ 的无向面积元为
\[
dS=\sqrt{1+z_x^2+z_y^2}\,dx\,dy.
\]
这条图形公式隐含一个表示资格：当前曲面片必须能在所选投影平面上写成单值图形。若投影退化成曲线（例如圆柱侧面投到与母线垂直的圆盘边界）或同一投影点对应多个曲面点，就应改投到其他坐标平面、分片，或直接参数化；不能对零面积/多重覆盖的投影硬套图形公式。曲面方程上的恒等式可以化简 $f$，但不能把 $dS$ 改写成 $dx\,dy$ 而漏掉斜率因子。

\begin{examplebox}{先代曲面方程，再决定微元}
球面 $x^2+y^2+z^2=R^2$ 上，
\[
\iint_\Sigma(x^2+y^2+z^2)\,dS
=R^2\operatorname{Area}(\Sigma)=4\pi R^4.
\]
同一方程若出现在第二型积分中，系数仍可代入，但有向微元不能因此变成 $dS$。
\end{examplebox}

\begin{corebox}{为什么线/面上的等式能代入，而区域边界不能当等式代入？}
积分载体若是曲线或曲面本身，例如
\[
L:\ g(x,y)=0,
\qquad
\Sigma:\ F(x,y,z)=0,
\]
那么载体上的\textbf{每一个积分点}都满足这个等式，所以可以用它化简被积函数；这只是利用载体上的点态恒等式。

但二重/三重积分的区域通常写成
\[
D:\ g(x,y)\le0,
\qquad
\Omega:\ F(x,y,z)\le0.
\]
这里等号只描述边界，内部点一般不满足 $g=0$ 或 $F=0$，因此不能把边界等式强行代进体内被积函数。稳定判断不是“线面积分能代、重积分不能代”的口诀，而是：
\[
\boxed{\text{当前积分点是否逐点满足这个等式？}}
\]
如果答案为是，就可作为恒等约束化简；若它只负责圈定允许范围，就只能用于区域编码。
\end{corebox}

\section{第二型曲线：向量场乘有向位移}

设 $\mathbf F=(P,Q,R)$，
\[
\boxed{\int_L\mathbf F\cdot d\mathbf r
=\int_L P\,dx+Q\,dy+R\,dz.}
\]
参数化方向必须与题设方向一致：
\[
\int_L\mathbf F\cdot d\mathbf r
=\int_\alpha^\beta\mathbf F(\mathbf r(t))\cdot\mathbf r'(t)\,dt.
\]
反向曲线整体变号。参数化时 $x,y,z$ 与 $dx,dy,dz$ 必须同步替换；只代坐标而漏微分，是第二型曲线的首个典型错误。

若曲线按题设方向取单位切向量 $\boldsymbol\tau$，则
\[
d\mathbf r=\boldsymbol\tau\,ds,
\qquad
\boxed{\int_L\mathbf F\cdot d\mathbf r
=\int_L(\mathbf F\cdot\boldsymbol\tau)\,ds.}
\]
这不是另一套曲线积分理论，而是把向量场先投影到运动方向再按弧长累积。它在题目只给某个抽象函数连续、无法合法调用需要一阶偏导的 Green/Stokes 时尤其有用：先利用切向量与曲线上的点态恒等式，看抽象项是否在切向投影中消掉，再决定是否参数化。不能为了套积分定理擅自把“连续”升级为 $C^1$。

若 $\mathbf F=\nabla\phi$，则
\[
\boxed{\int_A^B\mathbf F\cdot d\mathbf r=\phi(B)-\phi(A)}
\]
与路径无关。判断无旋能否升级为保守，还需检查定义域是否单连通。

\section{第二型曲面：向量场乘有向面积}

\[
\boxed{\iint_\Sigma\mathbf F\cdot d\mathbf S
=\iint_\Sigma P\,dy\,dz+Q\,dz\,dx+R\,dx\,dy.}
\]
参数曲面定向由 $\mathbf r_u\times\mathbf r_v$ 决定：
\[
d\mathbf S=(\mathbf r_u\times\mathbf r_v)\,du\,dv.
\]
若 $\mathbf n$ 是题设侧的单位法向量，则同一关系也可写成
\[
\boxed{d\mathbf S=\mathbf n\,dS,
\qquad
\iint_\Sigma\mathbf F\cdot d\mathbf S
=\iint_\Sigma(\mathbf F\cdot\mathbf n)\,dS.}
\]
因此第二型曲面积分就是“先取向量场的法向分量，再按无向面积累积”；这与上一节 $d\mathbf r=\boldsymbol\tau ds$ 完全平行。所谓两类曲面积分的转换，不需要再背三组方向余弦口诀，核心就是有向面积向量等于单位法向量乘无向面积元。交换参数次序会反转叉积与通量符号。

若 $z=z(x,y)$，上侧非单位法向量为 $(-z_x,-z_y,1)$，故
\[
\boxed{
\iint_\Sigma\mathbf F\cdot d\mathbf S
=\iint_D\bigl(-Pz_x-Qz_y+R\bigr)\,dx\,dy.}
\]
下侧整体反号；封闭曲面则以外法向为正，内向结果加负号。

\begin{warnbox}{有向投影不是普通面积}
$dy\,dz,dz\,dx,dx\,dy$ 是带方向的投影面积微元。一般法向量不能直接与任意坐标平面微元相乘；必须使用与投影/参数化匹配的有向面积向量。
\end{warnbox}

\begin{corebox}{直觉锚点：场值和方向各管一件事}
在坐标面 $z=0$ 上，若向量场第三分量是 $R(x,y,z)=xy$，代入曲面后仍然有 $R(x,y,0)=xy$；“$z=0$”不会凭空给 $xy$ 乘一个零。若该面取下侧，改变的是有向面积向量：
\[
d\mathbf S=-\mathbf e_z\,dx\,dy,
\]
所以通量贡献是 $-xy\,dx\,dy$。

因此始终把两层分开：
\[
\boxed{\text{坐标代入决定场有多大}\qquad\text{法向/定向决定取哪个分量以及什么符号}.}
\]
最小检查法是先在脑中问：“这里流体本身的速度分量是多少？”再问：“这块面朝哪边？”两问不能合成一步。
\end{corebox}

\begin{corebox}{抽象函数只给连续：不要主动制造更强资格}
若题目把一个复杂项写成任意连续函数 $f(\circ)$，这往往是在提醒：答案不应依赖 $f$ 的具体形式，也不应通过对 $f$ 求偏导来完成。先把场拆成“抽象部分 + 显式部分”，与有向微元逐层点乘，再利用曲面方程检查抽象部分是否结构性消去。

这是一条很重要的路由原则：
\[
\boxed{\text{题目只给到哪一级正则性，就优先寻找只需要这一级正则性的路线}.}
\]
Gauss/Stokes 若需要 $C^1$，而题目只给连续，就不能为了使用漂亮定理擅自升级条件；直接投影、参数化、载体恒等式或对称性可能反而是更高阶的技巧。
\end{corebox}

\section{对称性：整体微元的符号先于奇偶口诀}

第一型积分只需检查区域/曲面和标量密度的对称；第二型积分还必须检查定向和微元在对称变换下的符号。最稳妥的做法是把参数化后的完整 integrand 当成一个整体变换，而不是只看 $P,Q,R$ 的表面奇偶。

关于原点的封闭曲面对第二型面积分，外法向也同时反向；因此“奇倍、偶零”必须以整体变换验证，不能把平面反射的对角/非对角表格直接移植。

\section{三大定理：内部微分量与边界积分}

三大定理共享同一结构：
\[
\boxed{\text{区域内部的微分量累积}=\text{区域边界上的积分}.}
\]

\subsection{Green：平面区域与闭曲线}

若 $D$ 有界、边界分片光滑且取正向（外边界逆时针、孔边界顺时针），$P,Q\in C^1$，则
\[
\boxed{\oint_{\partial D}P\,dx+Q\,dy
=\iint_D(Q_x-P_y)\,dA,}
\]
\[
\boxed{\oint_{\partial D}P\,dy-Q\,dx
=\iint_D(P_x+Q_y)\,dA.}
\]
第一式是环流/二维旋度，第二式是平面通量/二维散度。非闭曲线可补线；有孔或奇点必须挖洞并保留内边界方向。

Green 还可以反向把面积编码成边界积分。对正向边界 $L=\partial D$，取 $P=-y/2,Q=x/2$，因为 $Q_x-P_y=1$，故
\[
\boxed{
\operatorname{Area}(D)
=\iint_D1\,dA
=\frac12\oint_L(x\,dy-y\,dx).}
\]
因此“区域内部面积难直接算、但边界易参数化”时，可以把面积问题改写为有向边界积分；若边界取负向，结果整体变号，所以几何面积最后必须保持非负并与方向协议一致。

\subsection{Stokes：定向曲面与空间边界曲线}

若 $L=\partial\Sigma$，曲线方向与曲面法向满足右手定则，且场在包含曲面的邻域内 $C^1$，则
\[
\boxed{\oint_L\mathbf F\cdot d\mathbf r
=\iint_\Sigma(\nabla\times\mathbf F)\cdot d\mathbf S.}
\]
同一边界的两张曲面在资格满足时可以互换；应比较替代曲面的投影域、旋度点乘、方向和分段总成本，而不是固定“优先投影”。曲线不是完整边界、定向不协调或场穿过奇点时不能直接使用。

\subsection{Gauss：体区域与封闭曲面}

若 $\Sigma=\partial\Omega$ 为封闭分片光滑曲面，取外法向，$\mathbf F\in C^1$ 覆盖整个体区域，则
\[
\boxed{\iint_{\partial\Omega}\mathbf F\cdot d\mathbf S
=\iiint_\Omega\nabla\cdot\mathbf F\,dV.}
\]
开曲面可补面；内部奇点要挖去小球并计算补偿通量。$\nabla\cdot\mathbf F=0$ 只控制无奇点区域的封闭总通量，不代表任意开曲面通量为零。

\section{散度与旋度的局部意义：把“小边界上的总量”还原成点密度}

Gauss 与 Stokes 不只是计算公式；把区域缩到一点，它们还告诉我们散度与旋度为什么是\textbf{局部密度}。

设 $\mathbf F\in C^1$，点 $P$ 周围取一族直径趋零的小体积 $\Omega_\varepsilon$。由 Gauss，
\[
\frac{1}{\lvert\Omega_\varepsilon\rvert}
\iint_{\partial\Omega_\varepsilon}\mathbf F\cdot d\mathbf S
=
\frac{1}{\lvert\Omega_\varepsilon\rvert}
\iiint_{\Omega_\varepsilon}\nabla\cdot\mathbf F\,dV
\longrightarrow
\boxed{\nabla\cdot\mathbf F(P)}.
\]
所以散度是“单位体积的净外流”在无限小尺度下的极限：正值表示局部源，负值表示局部汇，零表示一阶尺度上没有净源汇。它不等于“场本身向外”或“任意开曲面通量为零”。

同理，取法向为单位向量 $\mathbf n$ 的小定向面片 $\Sigma_\varepsilon$，边界为 $L_\varepsilon=\partial\Sigma_\varepsilon$。由 Stokes，
\[
\frac{1}{\lvert\Sigma_\varepsilon\rvert}
\oint_{L_\varepsilon}\mathbf F\cdot d\mathbf r
=
\frac{1}{\lvert\Sigma_\varepsilon\rvert}
\iint_{\Sigma_\varepsilon}(\nabla\times\mathbf F)\cdot\mathbf n\,dS
\longrightarrow
\boxed{(\nabla\times\mathbf F(P))\cdot\mathbf n}.
\]
因此旋度沿 $\mathbf n$ 的分量，是“单位面积的局部环流密度”。改变 $\mathbf n$ 就是在探测旋度向量的不同分量；若所有方向上的小环流密度都为零，才得到 $\nabla\times\mathbf F(P)=\mathbf0$。

\begin{corebox}{“任意闭边界恒为零”为什么能反推出局部方程？}
若连续函数 $g$ 在区域 $U$ 中满足
\[
\iiint_\Omega g\,dV=0
\]
对 $U$ 内每个足够小的体区域 $\Omega$ 都成立，则必有 $g\equiv0$。否则若 $g(P)>0$（或 $<0$），连续性会保证 $P$ 附近一小块区域仍保持同号，积分不可能为零。

因此当题目说“对半空间/区域内\textbf{任意}光滑闭曲面 $\Sigma$，通量恒为零”时，正确压缩链是
\[
\boxed{
\text{任意闭曲面通量}=0
\xRightarrow{\mathrm{Gauss}}
\iiint_\Omega\nabla\cdot\mathbf F\,dV=0\ \text{对任意小 }\Omega
\Longrightarrow
\nabla\cdot\mathbf F=0.}
\]
若场中含未知一元函数，这一步往往把积分题直接压成 Topic12 可处理的常微分方程。一个固定闭曲面的积分为零绝对不足以推出散度处处为零；真正提供点态信息的是\textbf{区域的任意性 + 连续性}。
\end{corebox}

\section{算子复合：梯度、散度与旋度怎样接起来}

梯度把标量场变成向量场，散度再把向量场压回标量场。因此对 $u\in C^2$，
\[
\boxed{
\Delta u:=\nabla\cdot(\nabla u)
=u_{xx}+u_{yy}+u_{zz},}
\]
称为标量场的 Laplace 算子（拉普拉斯算子）。这里保留的是数学一会直接计算和组合的算子接口；一般偏微分方程理论不进入本册。

若 $u(\mathbf x)=f(\rho)$ 是三维径向函数，$\rho=\sqrt{x^2+y^2+z^2}>0$，则链式法则合并三项二阶偏导后得到
\[
\boxed{
\Delta f(\rho)=f''(\rho)+\frac{2}{\rho}f'(\rho)
=\frac1{\rho^2}\bigl(\rho^2f'(\rho)\bigr)'.}
\]
所以题面若给 $\nabla\cdot\nabla f(\rho)=0$，三元偏微分表达式会因径向对称性降成一元 ODE：
\[
(\rho^2f'(\rho))'=0
\quad\Longrightarrow\quad
f(\rho)=C_2-\frac{C_1}{\rho}
\qquad(\rho>0).
\]
这一步只负责“向量算子条件 $\to$ 一元 ODE”的翻译；一元方程求解机制由 Topic12 接收。\textbf{$\rho=0$ 必须单独审计：}$1/\rho$ 在原点奇异；若题目还要求函数在包含原点的区域内具有足够正则性，则通常会进一步排除该奇异支，不能把穿孔区域上的通解无条件延拓到原点。

另一方面，若 $u\in C^2$，混合偏导可交换，故
\[
\boxed{\nabla\times(\nabla u)=\mathbf0.}
\]
它说明光滑势函数生成的梯度场必然无旋；反向“无旋 $\Rightarrow$ 全局梯度场”仍需定义域拓扑资格，不能跳过前面的穿孔反例。

\begin{methodbox}{定理选择不是固定排名，而是“资格门 + 成本比较”}
边界是平面闭曲线，可把 Green 列为候选；空间闭曲线可列 Stokes；封闭曲面可列 Gauss。先过“定向、光滑性、定义域奇点、边界完整性”四道资格门，再比较转换后的区域、integrand、分段数和补偿项。\textbf{定理可用不等于定理一定更快}；若换过去更复杂，就保留直接参数化/投影。
\end{methodbox}

\section{贯穿母例：把复杂边界换成简单对象}

\begin{examplebox}{球面的外向通量}
对半径 $R$ 球面和 $\mathbf F=(x,y,z)$，外法向满足 $\widehat{\mathbf n}=(x,y,z)/R$，所以 $\mathbf F\cdot\widehat{\mathbf n}=R$，
\[
\iint_\Sigma\mathbf F\cdot d\mathbf S=4\pi R^3.
\]
也可用 Gauss：$\nabla\cdot\mathbf F=3$，体积为 $4\pi R^3/3$。
这里“场 $\times$ 载体 $\times$ 定向”依次固定为向量场、封闭球面和外法向；直接法与 Gauss 法给出同一结果，且量纲为长度三次。若把球面当成无向面积元，方向信息会在第一步丢失。
\end{examplebox}

\begin{examplebox}{穿孔区域阻断“无旋即保守”}
在 $\mathbb R^2\setminus\{0\}$ 上，
\[
\mathbf F=\left(-\frac y{x^2+y^2},\frac x{x^2+y^2}\right)
\]
局部旋度为零，但单位圆逆时针环流为 $2\pi$。这里最有用的直觉不是背结果，而是认出
\[
\mathbf F\cdot d\mathbf r
=\frac{-y\,dx+x\,dy}{x^2+y^2}=d\theta
\]
在任意能连续选取角度的局部区域都成立。沿闭路走完一圈，连续抬升的角度变化量是 $2\pi\,\operatorname{wind}(L,0)$，所以
\[
\boxed{\oint_L\mathbf F\cdot d\mathbf r=2\pi\,\operatorname{wind}(L,0).}
\]
因此“绕不绕奇点、绕几圈、方向怎样”才是控制量，路径具体形状不是。奇点和非单连通性共同阻断全局单值势函数；挖洞不是技术补丁，而是把被破坏的拓扑资格显式化。
\end{examplebox}

\section{物理量与结果检查}

第一型曲线/曲面积分可建模线质量、面质量、质心与形心；第二型曲线是变力做功/环流，第二型曲面是通量。检查时至少确认：
\begin{itemize}
  \item 第一型结果不随反向改变，第二型结果反向变号；
  \item 通量、做功和质量的量纲与密度/场的定义一致；
  \item 只有“几何载体 + 密度”共同保持某个对称时，质心才被相应对称元素固定；
  \item 闭曲线/曲面换表示后的结果与直接法一致。
\end{itemize}

\section{专题执行协议}
\begin{enumerate}
  \item \kw{Classify}：判断标量/向量、曲线/曲面、第一型/第二型；
  \item \kw{Orient}：标注曲线方向、曲面侧向、边界右手关系；
  \item \kw{Measure}：写清 $ds,dS,d\mathbf r,d\mathbf S$，不可跨类型替换；
  \item \kw{Qualify}：检查闭合、单连通、$C^1$、奇点和边界完整性；
  \item \kw{Choose}：比较直接参数化、投影、势函数、补线/补面和三大定理；
  \item \kw{Compute}：同步代入坐标和微元，按方向保留符号；
  \item \kw{Verify}：反向、对称、量纲、简单场和换面不变量复核。
\end{enumerate}

\section{高频失败路径：第一次偏离发生在哪里？}
\begin{itemize}
  \item 把 $ds,dS$ 与 $dx,dy,dz,d\mathbf S$ 混成同一种微元；
  \item 参数化只代坐标，不同步代换微分或起终点方向；
  \item 图形曲面通量漏掉上/下侧或外/内侧符号；
  \item 只看系数奇偶，不看定向和完整有向微元；
  \item 无旋场直接宣布保守，漏查穿孔区域；
  \item Green 公式外/内边界方向相同，或开曲线补线方向错误；
  \item Stokes 曲线不是完整边界、右手方向不协调或换曲面穿过奇点；
  \item Gauss 对开曲面直接使用，或忽略封闭体内奇点；
  \item 看到复杂公式就直接参数化，漏掉势函数、投影或三大定理；
  \item 把梯度、散度、旋度的局部解释升级为未经资格审计的全局结论。
\end{itemize}

\section{传统章节与 Source Corpus 映射}
\begin{longtable}{L{2.75cm}L{5.15cm}L{\dimexpr\linewidth-7.90cm\relax}}
\toprule
\textbf{Source / 传统内容} & \textbf{进入本册} & \textbf{分流与拒绝复制}\\
\midrule
II-05 第一型曲线/曲面 & 标量密度、无向微元、参数化、投影、质心 & 无向面积体积接口调用 Topic08\\
II-06 第二型曲线 & 切向位移、方向、势函数、Green/Stokes 选择 & 线代内积统一解释归 B00\\
II-07 第二型曲面 & 法向面积、投影、外/内侧、Gauss & 一般微分形式不进入主干\\
II-08 统一框架 & 四类积分矩阵、微元不可混换、方法决策 & 作为 Topic09 组织骨架\\
II-08.1 梯度/散度/旋度 & 局部场读法与微分恒等式 & PDE/流形拓展不进入主干\\
II-08.1 保守场/拓扑 & 单连通资格、穿孔反例、路径无关 & 一般同调理论只留边界说明\\
II-08.1 Green & 环流/通量、补线、挖洞、边界方向 & 二重积分计算调用 Topic08\\
II-08.1 Stokes & 旋度通量、换面不换边界、右手定则 & 曲面参数化机制保留接口\\
II-08.1 Gauss & 封闭通量、补面、挖洞与源汇 & 三重积分计算调用 Topic08\\
\bottomrule
\end{longtable}

\section{一页压缩：忘掉公式后怎样重建？}
\begin{corebox}{一句话}
先问“什么场乘什么微元”，再问“方向由谁规定”；能否换成边界定理，取决于闭合、正则、单连通和奇点资格。第一型看无向长度/面积，第二型看切向/法向分量，Green、Stokes、Gauss 都把内部微分量转换成边界积分。
\end{corebox}

\subsection{四个重建公式}
\[
\int_L f\,ds=\int f(\mathbf r(t))\lVert\mathbf r'(t)\rVert\,dt,
\qquad
\int_L\mathbf F\cdot d\mathbf r=\int\mathbf F(\mathbf r(t))\cdot\mathbf r'(t)\,dt,
\]
\[
\iint_\Sigma f\,dS=\iint f(\mathbf r)\lVert\mathbf r_u\times\mathbf r_v\rVert\,du\,dv,
\qquad
\iint_\Sigma\mathbf F\cdot d\mathbf S=\iint\mathbf F(\mathbf r)\cdot(\mathbf r_u\times\mathbf r_v)\,du\,dv.
\]
\[
\oint_{\partial D}\text{边界形式}=\iint_D\text{二维微分量},
\qquad
\oint_{\partial\Sigma}\mathbf F\cdot d\mathbf r=\iint_\Sigma(\nabla\times\mathbf F)\cdot d\mathbf S,
\]
\[
\iint_{\partial\Omega}\mathbf F\cdot d\mathbf S=\iiint_\Omega\nabla\cdot\mathbf F\,dV.
\]

\subsection{重建问题}
\begin{enumerate}
  \item 当前是标量密度还是向量场？几何载体是线还是面？
  \item 微元是无向长度/面积，还是有向位移/面积？
  \item 曲线/曲面方向、边界方向和法向是否满足同一右手协议？
  \item 能否用曲面方程直接化简系数，但仍保留正确微元？
  \item 是平面闭曲线、空间闭曲线还是封闭曲面？
  \item 场是否在相关区域 $C^1$，定义域是否单连通，内部是否有奇点？
  \item Green、Stokes、Gauss 中哪个能把边界换成更简单对象？
  \item 结果能否通过反向、对称、量纲和简单场复核？
  \item 是否已经越界到 Topic08、Topic07 或 B00？
\end{enumerate}

\begin{center}\rule{0.5\linewidth}{0.5pt}\end{center}
\appendix
\section{定向积分 Coverage 锚点}
\begin{center}
\begin{tabular}{L{2.55cm}L{5.15cm}L{\dimexpr\linewidth-7.70cm\relax}}
\toprule
层级/对象 & 核心资格 & 输出\\
\midrule
第一型曲线/曲面 & 标量场 + 无向 $ds/dS$ & 长度/面积/质量\\
第二型曲线 & 向量场 + 有向 $d\mathbf r$ & 做功/环流\\
第二型曲面 & 向量场 + 有向 $d\mathbf S$ & 通量\\
保守场 & 单值势函数/单连通资格 & 端点差\\
Green/Stokes & 闭边界、正则性、定向协调 & 二维旋度/旋度通量\\
Gauss & 封闭边界、外向、无内部奇点 & 体内散度总量\\
\bottomrule
\end{tabular}
\end{center}

本册覆盖四类曲线/曲面积分、参数化与投影、保守场、Green、Stokes、Gauss、定向、补线/补面/挖洞和结果审计；Topic08 接收二重/三重积分，Topic07 提供梯度局部模型，B00 提供内积/投影统一解释。

\end{document}
